\begin{theorem}[All four elementary physical route operations]
    \label{thm:peps_gauged_route_step}
    \lean{
        TNLean.PEPS.IsGIsometric.exists_unitary_torusRouteSteps,
        TNLean.PEPS.torusRouteStepRegion,
        TNLean.PEPS.torusRouteStepAssignment
    }
    \leanok
    Let both periods of a finite regular torus be at least four, and suppose the
    original site tensor is $G$-isometric for the regular virtual action. There is
    one family of six-spin unitaries, indexed by the actual horizontal or vertical
    movement tile and its starting vertex. For every tile, common vertex gauge,
    insertion label, boundary configuration and literal exterior and crossing
    assignment, the forward unitary extends the single-chord insertion to the
    second chord. Its adjoint reverses the movement. Thus the family supplies
    rightward, leftward, upward and downward elementary operations on the actual
    open contractions and the corresponding globally contracted states. Both
    periodic seams are included.

    This is a finite-torus common-background consequence of SCP10,
    Theorem~6.16 and the accessible-coordinate construction (local lines
    1765--1920 and 2271--2305). Applying the family to a complete endpoint route
    still requires the literal assignment at every intermediate step; no
    four-endpoint braid or reunion protocol is asserted here.
    See \cite{gap:rmp_peps_quantum_double_g_isometry}.
    \uses{thm:peps_gauged_horizontal_flux_move,thm:peps_torus_gauged_vertical_flux_move}
\end{theorem}
\begin{proof}\leanok
    Choose the horizontal and vertical witnesses at every starting vertex,
    before all insertion and boundary data. Select the witness by the movement
    axis; select the witness or its adjoint by the direction. The two previously
    proved common-background identities give the actual original-spin actions.
\end{proof}
